\documentclass[12pt]{article}
\usepackage[utf8]{inputenc}
\usepackage[T1]{fontenc}
\usepackage{amsmath,amssymb,tikz}
\usepackage{geometry}
\geometry{margin=2.5cm}

\begin{document}

\begin{center}
\textit{--4--}
\end{center}

\noindent\fbox{Proof of Prop.} \quad $\begin{pmatrix}\cos\alpha & -\sin\alpha\\ \sin\alpha & \cos\alpha\end{pmatrix}^{\!n} = \begin{pmatrix}\cos n\alpha & -\sin n\alpha\\ \sin n\alpha & \cos n\alpha\end{pmatrix}$, \quad ord $R_{\frac{2\pi}{n}} = n$

\bigskip
\noindent $S\neq I_2$ and $S^2=I_2 \Rightarrow$ ord $S=2$ \hfill direct verification

\bigskip
\bigskip
\noindent I show that $\langle R,S\rangle$ has $2n$ elements.

\noindent Let $H=\langle R\rangle = \{I_2,R,\ldots,R^{n-1}\}$

\noindent $S\notin H$ (otherwise $S=R^i$)

\noindent $\det S = -1 = (\det R)^i = 1^i = 1$ \quad $\times$)

\bigskip
\noindent $HS = \{R^iS \mid i=\overline{0,n-1}\} \Rightarrow H\cap HS=\varnothing$

\noindent Hence $H\cup HS \subseteq D_n \Rightarrow$ ord.\ $D_n\geq 2n$

\bigskip
\noindent\underline{Remark} \ If $G$ is a group, $H\subseteq G$, $H$ finite and stable

\noindent under multiplication ``$\cdot$'' $\Rightarrow H\leq G$.

\bigskip
\noindent\underline{Proof} \ Let $a\in H$ and consider $aH \overset{\text{def}}{=} \{ax \mid x\in H\}$

\noindent Then we have $aH=H$.

\noindent ``$\subseteq$'' evident.

\bigskip
\noindent But $\left.\begin{aligned} H&\to aH\\ x&\mapsto ax \end{aligned}\right\}$ injective\ $\Rightarrow aH=H$.

\bigskip
\noindent $\Rightarrow ax=a \Rightarrow x=e\in H.$

\noindent But $e=ab \Rightarrow ba=b \Rightarrow b=a^{-1}\in H.$

\bigskip
\noindent\underline{Remark.} that $H\cup HS \ni R,S \Rightarrow \langle R,S\rangle \subseteq H\cup HS$

\noindent $R^iS\cdot R^jS = R^{i-j}S$ etc.\footnote{this verification that $H\cup HS$ is closed under multiplication is very compressed in the original; the exact exponents could not be reconstructed with full certainty.}

\bigskip
\bigskip
\noindent\underline{Ex}: \quad $\Delta_6$

\begin{center}
\begin{tikzpicture}[scale=1.1]
\draw[->] (0,-2.3) -- (0,2.3) node[above] {$y$};
\draw[->] (-3,0) -- (3,0) node[right] {$x$};
\coordinate (v1) at (1.7,1);
\coordinate (v2) at (0,2);
\coordinate (v3) at (-1.7,1);
\coordinate (v4) at (-1.7,-1);
\coordinate (v5) at (0,-2);
\coordinate (v6) at (1.7,-1);
\draw (v1) -- (v2) -- (v3) -- (v4) -- (v5) -- (v6) -- cycle;
\node[above right] at (v1) {$1$};
\node[above] at (v2) {$2$};
\node[above left] at (v3) {$3$};
\node[below left] at (v4) {$4$};
\node[below] at (v5) {$5$};
\node[below right] at (v6) {$6$};
\draw[->] (0.9,0.6) arc (30:80:1.1);
\end{tikzpicture}
\end{center}

\noindent $R_{\frac{2\pi}{6}} = R$

\noindent $R^2, R^3,\ldots,R^5$

\end{document}
